% !TEX root = ../CINTA-cn.tex
%%%%%%%%%%%%%%%%%%%%% chapter02.tex %%%%%%%%%%%%%%%%%%%%%%%%%%%%%%
% Author: Libin Wang
% Chapter 3 Congruences and Powers
% This document was released as open source on 2026-06-13.
% Copyright 2023, Libin Wang. Creative Commons License
% This work is licensed under a Creative Commons Attribution-NonCommercial-NoDerivatives 4.0 International License.
%%%%%%%%%%%%%%%%%%%%%%%%%%%%%%%%%%%%%%%%%%%%%%%%%%%%%%%%%%%%%%%%%%
\chapter{同余与模运算}
\label{chap:cong} 

\begin{introduction}
  \item 同余
  \item 模算术运算
  \item 同余方程求解
  \item 模指数运算
\end{introduction}

\section{同余与模算术运算}
\label{sec:congruence}

本章重点关注计算机系统中整数的加、减、乘、除等算术运算。在开始学习之前用以下 C 程序作为热身，请读者先阅读程序，尝试给出程序的输出，然后编译运行它们，看看自己的答案是否对。答案与结论并不重要，只是邀请众位与数字玩耍，也许从中可以找到一些规律与乐趣。

\begin{center}
	\lstinputlisting[language = C, title = 热身练习-1.]{codes/chap1-1.c}
\end{center}

\begin{center}
	\lstinputlisting[language = C, title = 热身练习-2]{codes/chap1-2.c}
\end{center}

通过以上练习~\footnote{多年的课堂教学经验表明，许多计算机专业低年级同学在修学过程序设计课程后，对以上练习及相关概念依然颇陌生。值得警惕！}，想强调的是，在计算机系统中每一个数值都表达为定长的二进制形式，其中的所有运算都是\emph{定长计算}\index{fixed-size computation}。
定长的长度可能是$8$比特、$16$比特、$32$比特或者$64$比特，通常是$8$的倍数，但不大可能是$7$比特或者其他。
计算机系统最基础的计算是二进制数值的与、或、非、异或、移位等逻辑操作，其次就是\emph{模算术运算}，即二进制数值的加、减、乘、除等运算操作。 
在这一章我们将学习一种自然、方便的描述模算术运算的方法：\emph{同余}( congruence)\index{congruence} 。

如果整数$a$、$b$和$m$满足$m\mid (a - b)$，我们称\emph{$a$与$b$模$m$同余}，记为：
\begin{equation}
\label{eq:congruence}
a \equiv b \pmod{m}
\end{equation}
称整数$m$为同余式的\emph{模数}(\emph{modulus}) 。等价地，等式~\ref{eq:congruence}也意味着存在整数$k$使得：
\begin{equation}
\label{eq:congru}
b = a + km
\end{equation}
等式~\ref{eq:congru}让我们想起定理~\ref{thm:DivAlgo}(除法算法)，$a$除$m$的余数等于$b$ 除 $m$的余数，也就是说：
\begin{equation}
\label{eq:congru_eq}
a \equiv b \pmod{m} \iff (a \bmod m) = (b \bmod m)
\end{equation}
有理由相信，这两种不同的同余的表示方式具有同等的重要性，在不同的情形下选择正确的表示法会给我们带来方便。
%The two forms of congruence are equivalently important.

\begin{example} 
\[26 \equiv 8 \pmod{9} \qquad \text{且}\qquad 6 \equiv 55 \pmod{7}\]
因为
\[9 | (26 - 8)\qquad \text{且}\qquad 7 | (6 - 55),\]
或者，等价地：
\[ 8 = 26 - 2 \cdot 9 \qquad \text{且}\qquad 55 = 6 + 7 \cdot 7.\]
\end{example}

%% 20260520
在继续讲解模运算的各种属性之前，让我们稍微停顿一次，思考模运算的另一种相对隐含的直观内涵，即\emph{模运算是用一把尺子（模数）对特定的数值进行测量}，根据模运算测量出的结果把特定的数值分类放入不同的桶中。图~\ref{fig:cong_motivation}表示出这种有趣且意义不凡的直观~\footnote{该图片由 Codex CLI 辅助生成。}，不同的整数按照模$m$ 所得余数的多少而被放到了不同的桶中。要测量出信息，通常要求尺子的尺寸不能太小，以下统一要求模数$m \ge 2$。

\begin{thinking}
  当我们要求$m \ge 2$，就不能不提一下《具体数学~\cite{GKP94}第 $4$ 章的热身题：$a \equiv b \pmod 0$ 意味什么？留给读者思考。
\end{thinking}

%20260615
相信“根据模$m$ 所得余数”只是测量的一种。图~\ref{fig:cong_motivation}也清晰地展示出了公式~\ref{eq:congru_eq}的内涵，也可视为一种测量。即通过判断$a - b$ 是否会落在 $0$ 号桶 从而判断两个整数 $a$ 和 $b$ 是否落在同一个桶中。为了方便与未来的内容进行对比讨论，值得在这里写下这个属性。下一节将会展示其他的“测量”方法。

\begin{property}{($a$与$b$模$m$同余)}
\label{prop:congruenc}
\[ a \equiv b \pmod m \iff (a - b) \equiv 0 \pmod m\]
\end{property}

图~\ref{fig:cong_motivation}还说明，通过模运算和同余关系还可以得到关于整数的若干属性。例如，对任意正整数$a$，$a \bmod 2$把所有的数分为两类，展示的是$a$的奇偶性，模数$2$就是一把尺子。如果模数是 $m=4$，能测量出更多的信息吗？比如，对所有的奇数$a$，此时会有两种情形$a \equiv 1 \pmod{m}$和$a \equiv 3 \pmod{m}$，即把奇数进一步分成了属性不同的两个集合。在此意义上，确实可以说更大的模数能测出更多的信息。

%%% The figure was created by Codex. 20260615
  \begin{figure}[htbp]
    \centering
    \includegraphics[width=0.9\textwidth]{chapt3-1-mod-congruence-buckets.png}
    \caption{将同余理解为模 $m$ 的测量}
    \label{fig:cong_motivation}
  \end{figure}

\begin{thinking}
\label{thinking:cong}
	如果模数 $m = 8$ 会怎样？请做简单的计算，看所有奇数在模$8$时会被分为多少类，并从中体会出：模数越大，同余关系施加的约束就越紧。注意，这里只是特定地指类似$2$、$4$和$8$这样，前一个数是后一个数的因子时产生的约束。如果数与数之间互素，则体现不出这种属性。
\end{thinking}

在本章，这种直观体现在同余方程的求解，特别是模数与系数不互素的情况下。在以后讲中国剩余定理时，这种直观会作用在求解同余方程组当模数之间不互素的情形。考虑更广阔更抽象的范畴，抽象代数中商群的定义中，一个群“模”一个正规子群的概念也有类似的内涵。换而言之，模运算并不仅仅是整数专有的运算。读者此时如果觉得困惑完全可以不需理会以上讨论，因为在深入到这些内容之前，我们会回到相对简单的整数模运算。

模加法、模乘法可以最直接的方式定义如下，并写成以下命题的形式：
\begin{proposition}{}{congru}
%\label{pro:congru}
如果$a_1 \equiv b_1 \pmod{m}$而且$a_2 \equiv b_2  \pmod{m}$，则
\[a_1 \pm a_2 \equiv b_1 \pm b_2 \pmod{m} \qquad \text{且} \qquad 
     a_1  a_2 \equiv b_1 b_2  \pmod{m}\text{。}\]
\end{proposition}

\begin{proof}
容易，作为课后练习。\qed
\end{proof}

命题~\ref{pro:congru}是说，如果对两个整数$a$和$b$进行模$m$加法或者乘法，可以首先算出$a' = (a \bmod m)$ 和 $b' = (b \bmod m)$，然后
对 $a'$ 和 $b'$ 做模 $m$ 的加法或者乘法，得到最终答案。在某种意义上，同余的值将被视为同一个值。并且，通过使用这种方法可以赢得效率，因为我们总是在尽可能小的数值上进行计算，看以下实例。

\begin{example}
因为 $10000 \equiv 1 \pmod{3}$ ，且$998 \equiv 2 \pmod{3}$，所以：
\[ 10000 \cdot 998 \equiv 2 \pmod{3}\]
\end{example}

\begin{example}
口算$7^{500} \bmod 8$ 。因为$7 \equiv -1 \pmod{8}$，所以$7^{500} \equiv 1 \pmod{8} $。
\end{example}

类似地，可以口算$998^{2} \bmod 999 = 1$，因为$998 \equiv -1 \pmod{999}$。

关于同余关系，还有以下重要的概念与命题。
\begin{proposition}{}{congr_eq_rel}
%\label{pro:congr_eq_rel}
同余关系是一种等价关系。
\end{proposition}
\begin{proof}
容易，请读者自行完成。\qed
\end{proof}

%To write something about congruence class.
如果有一种等价关系定义在集合$\mathbb{S}$ 上，则该等价关系会把集合 $\mathbb{S}$ 划分为不相交的子集，
这些子集我们称为\emph{等价类}\index{等价类}(\emph{equivalence classes})\index{equivalence classes}。集合中任意两个元素落在相同的等价类中
当且仅当这两个元素等价。反之，集合的一个划分，即集合分割为不相交的子集，定义了该集合上的一个等价关系。

考虑整数集合 $\mathbb{Z}$，设$m$为正整数，则模$m$的同余关系把$\mathbb{Z}$划分为$m$个等价类。任意整数$a$所落在的等价类称为
$(a \bmod m)$的\emph{同余类}(\emph{congruence class})\index{congruence class} ，或者称为\emph{剩余类}(\emph{residue class})\index{residue class}，记为$[a]_m$，即所有模$m$同余的整数形成的集合，或者所有具有以下形式的整数的集合：
$a + k m$，$k \in \mathbb{Z}$ 。也可以把同余类理解为在同一把尺子（模数）的度量下，将尺寸相同的元素放在同一个集合，图~\ref{fig:cong_motivation}已经表述了这种思想。%20260520

模$m$同余类的集合记为$\mathbb{Z}/m\mathbb{Z}$~\footnote{一种将在第~\ref{chap:iso_hom}章才解释的表示法，目前只需要把其看为一个记号。}，该集合恰好有$m$个元素，即：
\[\mathbb{Z}/m\mathbb{Z} = \{[0]_m, [1]_m, \cdots, [m-1]_m\}.\]

\begin{example}
\label{exam:residue-class}
如果$m = 2$，$\mathbb{Z}/2\mathbb{Z} = \{[0]_2, [1]_2\}$。同余类$[1]_2$就是所有模$2$余$1$的整数，即奇数，而$[0]_2$则是所有偶数的集合。同理，如果$m = 3$，$\mathbb{Z}/3\mathbb{Z} = \{[0]_3, [1]_3, [2]_3\}$。
\end{example}

称同余类$[a]_m$中的任意一个元素$b$为\emph{代表元}(\emph{representative})\index{representative}。
通常，使用一个同余类中最小非负元素作为代表元来表示一个同余类，但是在特定情况下选用某些特定的元素也许会更方便。
$\Z/m\Z$ 的所有最小非负代表元的集合$\{0, 1, 2,\cdots, m - 1 \}$被称为模$m$的\emph{最小剩余系}(\emph{least residue system}) \index{least residue system}。任意$m$个整数的集合，如果其中任意两个元素都不模$m$同余，则称之为模$m$的\emph{完全剩余系}(\emph{complete residue system})\index{complete residue system}。

\begin{example}
设$m = 7$，模$m$的最小剩余系就是集合 
\[\{0, 1, 2, 3, 4, 5, 6 \}\text{，}\]
下面这个集合则是模$m$的一个完全剩余系： 
\[\{14, 8, 23, 10, 46, 61, 13\}\text{。}\]
\end{example}

使用同余类的表达，命题~\ref{pro:congru}可简洁地重写为如下形式：
\begin{proposition}{}{residue}
%\label{prop_residue}
如果$[a_1]_m = [a_2]_m$ 且 $[b_1]_m = [b_2]_m$，则
\[[a_1 \pm b_1]_m = [a_2 \pm b_2]_m \qquad \text{且}\qquad
[a_1 b_1]_m = [a_2 b_2]_m \text{。}\]
\end{proposition}

\begin{proof}
把$[a]_m = [b]_m$变形为$a \equiv b \pmod{m}$，然后使用命题~\ref{pro:congru}，立得。 \qed
\end{proof}

%20260620
\section{消去律与一元线性同余方程的求解}
这一节重点考虑模算术运算中的加法和乘法，以及一元线性方程的求解。首先，
利用命题~\ref{pro:congru}，我们可以解这样一种带一个未知变量的同余方程。
\begin{example}
求解以下带变量$x$的同余方程：
 \[x + 7 \equiv 3 \pmod{11}\]
因为\[ -7 \equiv -7 \pmod{11},\]
将上两式相加，得到：
 \[x \equiv 3 - 7 \equiv -4 \pmod{11}.\]
\end{example}

因此，有一个好消息：
	\begin{center}
	\begin{tcolorbox}[colback=green!5!white, colframe=red!75!black,width = 10cm]
对形如以下的同余式：
\[ x + a \equiv b \pmod{m}\text{，}\]
其解是：
\[ x \equiv b - a \pmod{m}\text{。}\]
	\end{tcolorbox}
	\end{center}
然而坏消息就是，要解以下带乘法的同余方程就有点困难： 
\[ ax \equiv b \pmod{m}\text{。}\]
也许我们希望可以像中小学生解方程一样，两边同“除”一个$a$：
\[ x \equiv \frac{b}{a} \pmod{m}\text{。}\]
一个直接的问题就是：$(\frac{b}{a} \bmod m)$可能并不存在。
到底什么是这里的除法？幸运的是，我们还有以下规则。

\begin{proposition}{消去律}{cancel}\index{Cancellation Law}
%\label{pro:cancel}
设$a, b, c \in \Z$，$m$是一个正整数。如果$\gcd(c, m) = 1$，且$ac \equiv bc \pmod{m}$，则$a \equiv b \pmod{m}$。
\end{proposition}

\begin{proof} 
	根据同余的定义和条件，有$m \mid (ac - bc)$，等价地，$m \mid (a - b)c$。因为$\gcd(c,m) = 1$，根据推论~\ref{cor:euclid11}，则$m \mid (a - b)$，即得。\qed
\end{proof}
以上证明巧妙地利用了同余的表示法，值得初学者重视。然后，必须问
消去律为什么要求 $c$ 与 $m$ 互素？借助图~\ref{fig:cong_motivation}可以把以上证明过程解释为：因为两个数 $ac$ 和 $bc$ 落在同一个桶中，所以 $ac - bc$ 会落在$0$号桶，即$m \mid (a - b)c$。然后，为了确保 $(a - b)$ 也落在$0$号桶，就只能要$c$不会为 $(a - b)$贡献 $m$ 中大于$1$的因子。

%%% 20251016
\begin{thinking}
	思考在中小学中学到的消去律：设$a, b, c$为正整数，如果$ac = bc$，则$a = b$。证明过程可以是：由条件得到$ac - bc = (a - b)c = 0$，因为$c\neq 0$，所以$a-b = 0$，即$a = b$。或者，也可以这样证明：因为 $ac = bc$，且$c \neq 0$，所以等式两边同除$c$，得到$a = b$。请问以上两种证明有什么区别？你更倾向于哪一种证明呢，为什么？
\end{thinking}

%%% 20260615, add a label which will be used later in the chapter 7
\begin{thinking}
\label{thk:ab_mod_m}
	设$a, b, m$为正整数，且$a b \equiv 0 \pmod{m}$。证明：$b \equiv 0 \pmod{m/d}$，其中$d =\gcd(a, m)$。
\end{thinking}

%求解一般线性同余方程
消去律只是告诉我们，方程$ac \equiv bc \pmod{m}$必然会导致$a \equiv b \pmod{m}$，也就是说可以在同余方程两边消去$c$。但是，在实际的方程求解中，我们还需要一个“除法”去“除”掉那个$c$。表面上看我们需要的是一个除法，但是实际上用的是乘法。
求解过程是，先求出一个记为$c^{-1}$的值，
它满足\[cc^{-1} \equiv 1 \pmod{m}\] 
然后将它乘在同余式的两边：
\[acc^{-1} \equiv bcc^{-1}\pmod{m},\]
\noindent 因此得到：
\[a \equiv b \pmod{m}.\]
	
这个$c^{-1}$被称为$c$在模$m$下的\emph{乘法逆元}(multiplicative inverse) \index{multiplicative inverse}\index{inverse}。注意，既然强调
是\emph{模 $m$ 下}的乘法逆元，就只考虑小于 $m$ 的 $c^{-1}$，或者所有模 $m$ 同余的逆元都视为同一个逆元。在不引起误会的场合，为了
表述方便，往往会省略\emph{模$m$}这个限定语。

\begin{example}
设$m = 11$，$c = 3$，那么$c$模$m$的乘法逆元$c^{-1}$是$4$，因为$3 \cdot 4 \bmod m =  1$。虽然，$15$也满足$3 \cdot 15 \bmod m = 1$，
但是$15 \equiv 4 \pmod{m}$，所以，根据同余理论，它们被认为是同一个值。
\end{example}

给定互素的$c$和$m$，如何求乘法逆元$c^{-1}$？此时，上一章讲解的B\'{e}zout定理(定理~\ref{thm:bezout})重新登场。因为$\gcd(c,m)=1$,所以，存在$x,y\in \Z$使得：
\[cx + my = 1\text{，}\]
等式两边模$m$，得到$c x \equiv 1 \pmod{m}$，那么这个$x$就是我们所要求的$c^{-1}$。

以上分析说明，给定互素的$c$和$m$，$c^{-1}$必然存在。进一步$c^{-1}$会必然唯一吗？不要忘记，B\'{e}zout系数并不唯一。唯一性的具体证明留给读者作进一步探讨。
至此我们已经部分地解决了同余方程$a x \equiv b \pmod{m}$的求解问题，形成以下命题。

\begin{proposition}{求解一元线性同余方程--简化版}{cong_eq_simple}
	设$a, b, m \in \Z$，$m \ge 2$，且$\gcd(a,m)=1$，则一元线性同余方程
  \[ax \equiv b \pmod{m}\]
  在模$m$的意义下有唯一解，且该解为
  \[x \equiv a^{-1}b \pmod{m}\text{，}\]
  其中，$a^{-1}$表示模$m$下$a$的乘法逆元。
\end{proposition}
%% 20220912 
%% 20240314

\begin{example}
求解 $3 x \equiv 2 \pmod{11}$。首先，使用$\egcd$算法(算法~\ref{alg:egcd})计算模$11$下$3^{-1} = 4$。
具体而言，是算出B\'{e}zout系数$s$和$t$使得
\[ 3s + 11t = 1\]
此时，$s = 4$，$t = -1 $，然后对以上等式两边同时模整数$11$，则可得：
\[ 3 \cdot 4 \equiv 1 \pmod{11}\]
即，$3$模$11$的乘法逆元是$4$。因此，对所需求解的同余方程两边
同乘$3$的乘法逆元$4$，得到$x \equiv 8 \pmod{11}$。
\end{example}

\begin{thinking}
  以上实例中提示使用$\egcd$算法求乘法逆元。请读者思考，求乘法逆元的算法与$\egcd$算法有什么区别？这是课后作业。
\end{thinking}

\begin{thinking} 
根据B\'{e}zout定理(定理~\ref{thm:bezout})思考乘法逆元的唯一性。
已知B\'{e}zout系数不唯一，为什么相应的乘法逆元唯一呢？这里需要一些思考和计算，或者参考FINT~\cite{silverman2014friendly}的第六章，或者~\href{https://github.com/lbwang/CINTA-cn/blob/master/Notes/uniqueness-bezout.tex}{本书的授课手记}。
\end{thinking} 

使用命题~\ref{pro:cong_eq_simple}求解的同余方程$ax \equiv b \pmod{m}$有很大的局限性，要求$\gcd(a,m) = 1$。如果$a$与$m$不互素会如何？以下讨论要求$1 < \gcd(a,m) < m$，旨在提示读者该如何自行分析求解该问题，不给出完整的证明和求解过程。

%Motivation about cancellation.
在求解这一类同余方程前，先借助图~\ref{fig:cong_motivation}的直观，换一个视角去考察命题~\ref{pro:cong_eq_simple}的内涵。表面上看，命题~\ref{pro:cong_eq_simple}表述的是，给定$a, b, m$，$\gcd(a,m) = 1$，如何求解$x$。现在这样考虑：给定$a$和$b$两个桶，是否存在一个$x$值，使得$ax$桶必然与$b$在同一个桶？必然存在，分析如下。

把$x$理解为一个跑遍所有剩余类$[0]_m, [1]_m, \ldots, [m-1]_m$的变量，用给定的$a$去“乘”这些桶，会得到这些桶的一次置换, 得到一个完全剩余系，
\[[a\cdot 0]_m, [a\cdot 1]_m, [a\cdot 2]_m, \ldots, [a\cdot(m-1)]_m \text{。}\]

用数学的语言表示，即定义映射
\[
	\phi_a\colon [0..m-1] \to [0..m-1] \text{，}\qquad
	x \mapsto ax \bmod m \text{。}
\]
$\phi_a$ 把任意的 $x \in [0..m-1]$ 映射为$(ax \bmod m) \in [0..m-1]$。然后只需要证明$\phi_a$是一个单射，然而这是容易的，因为给定两个不同的值$x_1, x_2 \in [0..m-1]$，$\phi_a(x_1) \ne \phi_a(x_2)$，
否则有
\[a x_1 \equiv a x_2 \pmod m\]
利用消去律，立得$x_1 \equiv x_2 \pmod m$。矛盾！这种方法以后章节还会使用。

用这个视角去理解命题~\ref{pro:cong_eq_simple}之所以有意思，是因为当$a$与$m$互素，所有的桶乘以$a$只会重新排列模所有的桶，
而不会合并桶或者丢失某些桶。我们可以认为，桶乘以$a$之后不会损失模$m$的测量信息。但是，如果当$a$与$m$不互素情况就不同了。先看以下实例数据。

\begin{example}{}{}
考虑之前定义的映射$\phi_a$。对不同的合数 $m$，选择与 $m$ 不互素的 $a$，有
\[
\begin{aligned}
	m=6,\ a=2:\quad&
	([0]_6,[1]_6,\ldots,[5]_6)\\
	&\longmapsto
	([0]_6,[2]_6,[4]_6,[0]_6,[2]_6,[4]_6);\\[4pt]
	m=8,\ a=4:\quad&
	([0]_8,[1]_8,\ldots,[7]_8)\\
	&\longmapsto
	([0]_8,[4]_8,[0]_8,[4]_8,
	  [0]_8,[4]_8,[0]_8,[4]_8);\\[4pt]
	m=12,\ a=8:\quad&
	([0]_{12},[1]_{12},\ldots,[11]_{12})\\
	&\longmapsto
	([0]_{12},[8]_{12},[4]_{12},
	  [0]_{12},[8]_{12},[4]_{12},
	  [0]_{12},[8]_{12},[4]_{12},
	  [0]_{12},[8]_{12},[4]_{12}).
\end{aligned}
\]
这些映射会把不同的剩余类映射到了同一个剩余类，因此不是单射；
同时，一些剩余类没有原像，因此也不是满射。读者能发现这些映射的像呈现什么规律吗？
\end{example}

使用“模运算是测量”的视角继续以下的分析。有一个非常重要的关注点得指出，要求解方程$a x \equiv b \pmod{m}$，此时模数$m$测量的主体是什么？是$a$吗？不是，是$ax$。在这个指导思想下，
B\'{e}zout定理(定理~\ref{thm:bezout})再次闪亮登场。因为$\gcd(a,m) = d \ge 2$，B\'{e}zout定理说，存在$x_0, y_0\in \Z$使得：
	\[d = \gcd(a, m) = ax_0 + my_0\text{。}\]
等式两边同模$m$，得到：
\[ ax_0 \equiv d \pmod m \] 
并且，根据第~\ref{chap:euclid}的属性~\ref{prop:gcd1}，这个$a x_0$是模$m$可测量出来的最小正整数。

进一步，第~\ref{chap:euclid}章属性~\ref{prop:gcd2}表明“$a$ 与 $m$ 的所有线性组合都是 $d$ 的倍数，且所有$d$ 的倍数都是$a$ 与 $m$的整数线性组合”，也就是说，对任意的$x \in [0..m-1]$，用$m$去模$a x$都得到$d$的某个倍数。在模 $m$ 的意义下，有多少个不同的 $d$ 的倍数？$m/d$个，分别是：
\[ 0, d, 2d, \ldots, (m/d - 1) d \text{。}\]
也就是说，$a x \bmod m$只能落在这样剩余系之中（把$(m/d - 1) d $简写为 $m - d$ ）：
\[ [0]_m, [d]_m, [2d]_m, \ldots, [m-d]_m \text{。}\]

模$n$的最小完全剩余系中 $m$ 个同余类乘以 $a$ 后，同余类的个数缩减为 $m/d$个。因此我们说，模 $m$ 的尺子损失了 $d$ 倍的分辨率。接下来，考虑$a x \equiv b \pmod{m}$的求解，立即需要问：
\begin{thinking}
%\label{thk:cong_eq}
\begin{itemize}
\item 什么情况下同余方程有解或无解？

\item 如果同余方程有解，会有多少个模$m$不同余的解？
\end{itemize}
\end{thinking}

作为最后的提示，下面给出一个用得上的命题。最后，建议读者继续自主完成方程的求解分析，留作课后练习。

\begin{proposition}{同余式约分律}{congruence_reduction}
设 $a,b\in\Z$，$m,d$为正整数，则
\[
	ad\equiv bd\pmod{md}
	\iff
	a\equiv b\pmod m.
\]
\end{proposition}

\begin{proof}
根据同余的定义，
\[
\begin{aligned}
ad\equiv bd\pmod{md}
&\iff md\mid(ad-bd)\\
&\iff md\mid d(a-b).
\end{aligned}
\]
因为 $d\neq0$，所以
\[
md\mid d(a-b)
\iff
m\mid(a-b).
\]
再次根据同余的定义，
\[
m\mid(a-b)
\iff
a\equiv b\pmod m.
\]
因此
\[
ad\equiv bd\pmod{md}
\iff
a\equiv b\pmod m.
\]
\qed
\end{proof}

作为提示，最后给出求解一般线性同余式的方法如下。
\begin{proposition}{求解一元线性同余方程--广义版}{cong_eq_complete}
	设$a, b, m \in \Z$，$m \ge 2$，且$\gcd(a,m)=d$，
  \begin{itemize}
    \item 如果$d \nmid b$，则一元线性同余方程$ax \equiv b \pmod m$无解。
    \item 如果$d \mid b$，则一元线性同余方程$ax \equiv b \pmod m$有且仅有$d$个非同余解。整个解集为：
    \[ x\equiv x_0 + k\frac{m}{d}\pmod m, \qquad k=0,1,\ldots,d-1 \text{。}\]
    其中$x_0$可以通过使用命题~\ref{pro:cong_eq_simple}的方法求解方程$a'x \equiv b' \pmod{m'}$得到，其中$a' = a / d$，$b' = b / d$，$m' = m / d$。
  \end{itemize}
\end{proposition}

本节的最后补充一个简单的小命题，考虑模素数的乘法逆元。该内容可作为对\emph{非线性同余方程}（\emph{Nonliearn congruences}）介绍的开始，也算是再次讲解同余关系的应用。
\begin{proposition}{}{inv_mod_prime}
%\label{prop:inv_mod_prime}
设$p$为素数，且$a\in \N$。则$a^2 \equiv 1 \pmod{p}$ 当且仅当$a \equiv 1 \pmod{p}$ 或 $a \equiv -1 \pmod{p}$。
\end{proposition}

\begin{proof}
     $\Leftarrow$的证明简单\footnote{本书不使用常见的“充分性”和“必要性”的表述，取而代之，用$\Rightarrow$表示从左到右的推导，而$\Leftarrow$表示从右到左的推导。这样看似不准确的表述就是为了避免表述上的不准确，并带来理解的方便。请注意，A当且仅当B，则它们互为充分、必要条件。所以，单纯表述为“充分性”或者“必要性”都不准确，而应该表述为“A是B的充分条件”或者“A是B的必要条件”等，这就显得有点繁琐。另外，也避免在中文语境下对“A当且仅当B”的理解分歧。某些教材把“A当且仅当B”分为“A当B”和“A仅当B”，前者为“必要性”，后者为“充分性”。作者对此表述不作评价，只是不推荐使用这种表述方法。}。
	对于$\Rightarrow$，如果 $a^2 \equiv 1 \pmod{p}$则$p \vert (a^2  - 1)$，这意味$p \mid (a - 1)$或 $p \mid (a + 1)$，
等价于$a \equiv 1 \pmod{p}$或$a \equiv -1 \pmod{p}$。\qed
\end{proof}
命题~\ref{pro:inv_mod_prime}，本质上是提醒我们，在模数为素数的情形下，什么数的乘法逆元是自己本身。或者，不准确地说，
在模素数的情形下，对$1$开根号，得到一个$1$和一个$-1$。难道如果模数不是素数会有什么不同？以后再说。感兴趣的读者可提前考虑对$-1$开根号的情形，即什么数会满足$a^2 \equiv -1 \pmod{p}$ 。

\section{二进制补码}
\label{sec:two_comp}
本节重点考虑模算术运算中的减法。对于整数$a$, $b$和$m$，以下等式：
\[ a - b \bmod m\]
等价于
\[a + (-b) \bmod m\]
其中，$-b$为一个\emph{带符号数}(\emph{signed number})\index{带符号数}\index{signed number}. 

计算机中表示带符号数的最常用方法就是\emph{二进制补码}(\emph{two's complement})\index{二进制补码}\index{two's complement}。
使用模算术运算，可以清晰地解释二进制补码的设计原理。 

先回顾一下二进制补码的常规表述。假定需要表达一个$n$比特的带符号数，最高位比特称为符号位，用于表示正或负，其余$n-1$个比特用于表达数值。
$n$比特带符号数的范围是$[-2^{n-1}, 2^{n-1} - 1]$，其中非负整数的表达范围是$[0, 2^{n-1} -1]$，而负整数的范围则是$[-2^{n-1}, - 1]$。
相关表达规则如下：

\begin{itemize}
	\item 非负整数按常规二进制数表达存储在$n-1$个低位比特中，其符号位置为$0$。
	
	\item 负整数的表达分为三个步骤。设$1 \leq x \leq 2^{n-1}$，要表达负整数$-x$：首先构造正整数$x$的二进制补码表达（规则如上），
	然后对$x$的所有比特进行一次比特翻转（$0$比特翻转为$1$；$1$比特翻转为$0$），最后对所得的$n$比特加$1$（一次二进制加法）。
\end{itemize}

\begin{example}
	设带符号数的比特长度$n = 8$，计算$-1$和$-128$的二进制补码。

	\[
	\begin{array}{r@{\;}c@{\;}l}
	+1 & = & 00000001 \\
	\text{比特翻转} & & 11111110\\
	\text{加 }1 & & 00000001\\
	-1 & = & 11111111
	\end{array}
	\]

	\[
	\begin{array}{r@{\;}c@{\;}l}
	+128 & = & 10000000 \\
	\text{比特翻转} & & 01111111\\
	\text{加 }1 & & 00000001\\
	-128 & = & 10000000
	\end{array}
	\]
  以上计算表示 $ (+128) = (-128)$。实际上，在系统实现中，这个$+128$并不出现在 $8$ 比特正整数的范围内。
\end{example}

大部分刚接触二进制的初学者都会懂得这种“比特翻转再加$1$”的二进制补码操作，然而很少人会理解，负数为何如此表达。
常规的教科书会告诉读者，这是为了效率，为了避免出现$-0$等等。以下给出另一种解释二进制补码的方式。
对$n$比特的正整数$x$（$1 \leq x \leq 2^{n-1}-1$）取反所得的
负整数$-x$就是$ 2^n - x$。 
 
从目标出发，所谓求正整数$x$的负值，就是求得某个整数$x'$使得$x + x' = 0$，即：
\[x + x' \equiv 0 \pmod{2^n}\]
根据模算术运算的规则，既然$x<2^n$，那么$x' = 2^n - x$。又因为$2^n - x = \underbrace{111\cdots 11}_{n}  + 1  - x$，
且$\underbrace{111\cdots 11}_{n} - x$就等于对$x$的所有比特取反。所以，$x'$的值就等于对$x$的所有比特取反之后再加$1$。

有了二进制补码，模算术加法和模算术减法的实现就完全地统一起来了。换而言之，使用模加法和某些逻辑操作可实例化出模减法。
\begin{example}
计算$ (5 - 27) \bmod 2^8$。
\begin{itemize}
\item 分别写出$5$和$27$的二进制表达为：$0b00000101$和 $0b00011011$；

\item $-27$的二进制补码表达是$0b11100101$；检验可知$0b00011011 + 0b11100101 = 0b00000000$；

\item 则计算$5 + (-27)$就是计算$0b00000101 + 0b11100101 = 0b11101010$；检验可知$0b11101010$就是十进制的$-22$.
\end{itemize}  
\end{example}
考虑到，计算机内部的所有运算都是定长计算，$n$比特整数的运算必然是模$2^n$的运算，这是考虑表达负数时的前提条件。那么把负整数$-x$理解为$ 2^n - x$
就不能算是对二进制补码的“翻转加$1$”的解释，而是正本清源罢了。

\section{模指数运算}
\label{sec:mod-pow}
这一节关注模指数运算，一种非常重要的算术计算原语。所谓模指数运算，就是输入整数$x$、$y$ 和 $m$，要求 $y\geq 0$，$m\geq 2$，计算：
\[ x^y \bmod m \text{。}\]

SageMath 给出命令 \lstinline{pow} 来计算模指数，或者可以直接用 \lstinline{^} 和 \lstinline{%} 运算符。
以下是具体实例。
\begin{lstlisting}
sage: pow(2, 10, 1000)
24
sage: 2^10 % 1000
24
sage: pow(2, 1024, 3^100)
426135005318503491421173769021520445807198574681
\end{lstlisting}	
 
完成模指数运算最简单的方法就是重复地乘$x \bmod m$。但是，如果利用计算过程的中间数值，计算效率可以大大提高。该过程
重复对$x$进行平方，例子~\ref{exm:pow}展示了这种思路。

\begin{example}
\label{exm:pow}
求$2^{16} \bmod 11$。计算过程：
\[2^2 \bmod 11 = 4\] 
\[2^4 \bmod 11 = 4 \cdot 4 \bmod 11 = 5\]
\[2^8 \bmod 11 = 5 \cdot 5 \bmod 11 = 3 \]
\[2^{16} \bmod 11 = 3 \cdot 3 \bmod 11 = 9 \]
\end{example}

以上例子中的指数恰好是$2$的乘方，如果不是怎么办？比如，要求的是$2^{22} \bmod 11$。如下例所示。
\begin{example}
\label{exm:pow2}
求$2^{22} \bmod 11$。不难发现：
\[2^{22} = 2^2 \cdot 2^4 \cdot 2^{16} \]
所以 
\[2^{22} \bmod 11 = 4\cdot 5 \cdot 9  \bmod 11 = 4 \]
\end{example}

以上例子说明，可以通过不断计算某个数的平方，然后利用中间结果来求出最终的值，从而提高计算效率。但是，该计算过程如何表示为算法呢？考虑以下递归过程，
具体实现参见算法~\ref{alg:rec_pow}的描述。
由此可知，模指数运算的复杂性不是$\mathcal{O}(y)$次乘法，而是$\mathcal{O}(\log (y))$次乘法，效率得到极大提升。
\begin{equation}  
\label{eq:rec_pow}
x ^ y = 
\left\{  
\begin{array}{lr}  
(x ^{ \lfloor y/2 \rfloor })^2 &  \text{如果$y$是偶数；}\\  

x \cdot (x ^{ \lfloor y/2 \rfloor })^2 & \text{如果$y$是奇数。} \\  

\end{array}  
\right.  
\end{equation}

\begin{center}
\label{alg:rec_pow}
	\lstinputlisting[caption = 模指数运算--递归版本]{codes/chap2-3.py}
\end{center}

接下来展示如何把递归算法转换为迭代算法。
首先，把整数$y$视为多项式（或者一个二进制数串）：
\[ y = y_{n-1}2^{n-1} + y_{n-2}2^{n-2} + \ldots + y_1 2 + y_0 ,\] 
其中$y_i \in \{0, 1\}$。

其次，将$x^y$变形为：
\[x^y = \prod_{i=0}^{n-1}  x^{y_i2^{i}}\]

最后，从$i=0$开始不断计算$x^{2^{i}} \bmod m$，判断如果$y_i = 1$，则把该值乘进最终结果。注意，如果$y_i = 0$，$x^{y_i2^{i}}$等于$1$。
以下例子展示了其具体计算过程。

\begin{example}
	设$x=7$，$y= 10$，$m = 11$，
	计算$x^y \bmod m$。$y$的二进制表达为$1010$，因此计算如下：
	\begin{align*}
		y_0 &= 0, \quad x^{2^0} \equiv 7 \bmod m\\
		y_1 &= 1, \quad x^{2^1} \equiv 5 \bmod m\\
		y_2 &= 0, \quad x^{2^2} \equiv 3 \bmod m\\
        y_3 &= 1, \quad x^{2^3} \equiv 9 \bmod m
	\end{align*}
然后，将$y_i = 1$的项乘起来，得到$x^y = (5 \cdot 9) \bmod 11 = 1$.
\end{example}

该过程可描述为以下算法。
\begin{center}
\label{alg:ite_pow}
	\lstinputlisting[caption = 模指数运算--迭代版本]{codes/chap2-4.py}
\end{center}

该算法的正确性与效率显然可知。值得一提的是，如果读者已经写出了简单乘法的迭代版本，就会发现这两种算法的相似之处。有趣！模指数运算的优化改进算法则不在本书讨论范围。

%%% 20230424下午创建
\section*{章注}
\addcontentsline{toc}{section}{章注}
% 删除许多与教材语气不协调的句子，避免削弱正文的可信度。不过，我真的是被“震惊”到了。20260622
对计算机专业的大学生而言，模运算的重要性如何强调都不过份。如果读者对本章的热身练习感到陌生，则需要警惕。本章花了较多的篇幅重复强调模运算与同余关系的直观，希望在未来读者可以自觉地把这种“模”的思想运用到更抽象的代数结构中。
使用模运算去解释二进制补码显得非常自然，然而，奇怪的是，作者并没有发现有什么教科书这么去做，也许太简单而不值得出现在教材中。所以多年前作者把它写在了博客~\href{https://www.jianshu.com/p/502d236116ca}{《关于二进制补码的若干注记》}之中。模指数运算的重要性则通过许多教科书得到了体现，它与简单乘法算法
结构的相似就显得非常有趣。《算法导论》给出一种与本书描述有微妙差异的算法，为什么那种算法的效率更好，Cohen 的教材~\cite{cohen2013course}(第1章)给出了答案。

\begin{problemset}

  \item 编写一个 C 语言程序，用于验证在 C 语言中，无符号整数的加法、减法、乘法和除法都是模$2^n$的模运算，$n$根据变量定义不同可能是$8$、$32$或者$64$等。
  
  \item 编程计算若干整数及其十进制各位数字之和模 $3$ 的余数，猜测一个整数能被 $3$ 整除的判定规则，并使用同余关系证明该规则。%AI润色版本。
  
  \item 请证明命题~\ref{pro:congru}。
  
  \item 给定整数$p$和$N$，且$p \mid N$。证明：对任意整数$x$，
  \[ (x \bmod N) \bmod p = (x \bmod p)\]
  给出反例说明$(x \bmod p) \bmod N$ 并不一定等于 $(x \bmod N)$。
  %%% x % N = x - q N, (x % N) % p = (x- qN) % p = x % p
  
  %%% From Andre Weil, NT4B
  \item 设$m$是任意正奇数，请证明$(1^m + 2^m + \cdots + (m-1)^m) \equiv 0 \pmod{m} $ ~\cite{Wei79}。
  
  \item 对整数集合$\{1, 2, 3\}$可以定义多少种不同的等价关系？%源自AC0第1.1章 
  推而广之，对$n$个元素的集合可以定义多少种不同的等价关系？%Bell数 20250430
  
  \item 请证明同余关系是等价关系，即证明命题~\ref{pro:congr_eq_rel}。
  
  %2024年3月，该题出现在X，发现以下证明。
  %设 n = ab，一个合数，记m =2^a - 1，那么 2^a \equiv 1 \pmod{m}，所以(2^a)^b  \equiv 1 \pmod{m}，即m整除2^n - 1 。 
  \item 请证明，如果$n$是合数，则$2^n - 1$不会是素数。
  
  \item 手动计算以下模$m$下$a$的乘法逆元。(a) $m = 11$，$a = 5$；(b) $m = 121$，$a = 13$；(c) $m = 1021$，$a = 131$。 
  %ans：9, 28, 265
 
   \item 编写 C 语言程序完成模指数运算，即给定整数$x$，$y$和$m$为输入，计算并返回值$x^y \bmod m$。
 
   \item 定义 Fibonacci 数列如下：$F(0)=0$，$F(1)=1$，且对于$n \geq 2$，$F(n)=F(n-1)+F(n-2)$。所以，该数列是：$0, 1, 1, 2, 3, 5, 8, 13, 21, \dots$。如何能快速地求出$F(n)$呢？
   很幸运，我们有以下等式：
   $$\begin{bmatrix} 1 & 1\\ 1 & 0 \end{bmatrix}^n = \begin{bmatrix} F(n+1) & F(n)\\ F(n) & F(n-1) \end{bmatrix}$$
   虽然，看上去该算法需要一次矩阵的指数运算，但是借助快速指数运算的方法，这里可以产生一个快速求解$F(n)$的算法。请给出算法，并编程实现，C 语言或者 Python 都可以。
   %%% https://www.nayuki.io/page/fast-fibonacci-algorithms
   %%% https://www.geeksforgeeks.org/fast-doubling-method-to-find-the-nth-fibonacci-number/
 
 \item 给定互素的正整数$c$和$m$，请证明在$\bmod \; m$的意义上存在唯一确定的整数值$c^{-1}$，它使得$c c^{-1} \equiv 1 \pmod{m}$。
 
 \item 假设$x \equiv a \pmod{n}$，且$n = p \cdot q$，$p, q$是两个不相同的素数。请证明，因此必然有$x \equiv a \pmod{p}$和$x \equiv a \pmod{q}$。反之，如果有$x \equiv a \pmod{p}$和$x \equiv a \pmod{q}$，是否必然有$x \equiv a \pmod{n}$？请证明或证伪。
 %%% 简单题，为CRT准备。
 
 \item 思考另一个版本的消去律。设$a, b, c \in \Z$，$m$是一个正整数。如果$\gcd(c, m) = g$，$g$是大于一的整数，且$ac \equiv bc \pmod{m}$。
 请问此时能消去$c$得到某个同余式$a \equiv b \pmod{m'}$吗？这个$m'$会是什么？给出证明。% m' = m/g, where g=gcd(c,m)
 
 \item 求解以下同余方程：
  \begin{itemize}
  \item $3x \equiv 7 \pmod{17}$  %ans: x = 8 mod 17
  
  \item $6x \equiv 9 \pmod{23}$  %ans: x = 13 mod 23
  
  \item $6x \equiv 15 \pmod{21}$  %ans: x = 6, 13, 20 mod 21
  % 2024年10月布置作业并讲解.
  \end{itemize}
 
 \item 完成同余方程求解的一般性分析。设$a, b, m\in \Z$，求解$a x \equiv b \pmod{m}$，其中$g = \gcd(a,m)$。

 \begin{itemize}
 \item 同余方程在什么情况下无解？
 
 \item 如果同余方程有解，会有多少个模$m$不同余的解？
 \end{itemize}
  
 \item 编写一个 Python 程序计算乘法逆元，即输入互素的正整数$c$和$m$，返回$c^{-1}$，它使得$c c^{-1} \equiv 1 \pmod{m}$。要求：只返回为正整数的$c^{-1}$。
  
  \item 编写一个 Python 程序求解以下同余方程：
     \[ a x \equiv b \pmod{m}\text{。} \]
    即，给定整数$a$，$b$和$m$作为输入，返回所有满足方程的解$x$，或者给出无解告警。
   
  %%% 添加于2024年12月16日。  
  \item 二次探测(Quadratic Probing)法是 Hash 表解决冲突的常用策略之一。一种常见的Hash函数的定义如下：$H(k,i) \equiv H(k) + i^2 \pmod{p}$，其中$p$是一个素数，
  而$i\in \Z_{(p-1)/2}$。请分析证明为什么需要取一个素数$p$为模数。%%% 确保the first ⌊m/2⌋ probe sequences are distinct. https://www.cs.umd.edu/class/fall2020/cmsc420-0201/Lects/lect14-hashing.pdf
    
  \item 第一章的习题~\ref{prob:gf}定义了集合$\mathrm{GF}$及其元素的加法与乘法。一个已知结论：集合$\mathrm{GF}$的非零元素都有乘法逆元，即任取$a \in \mathrm{GF}$，且$a \ne 0$，必然存在$b \in \mathrm{GF}$使得$a\cdot b = 1$。请写一个程序\texttt{GF\_inverse}完成求逆元的任务，即对任意的$a \in \mathrm{GF}$，
    且 $a \ne 0$，求$b \in\mathrm{GF}$使得$a \cdot b = 1$。[提示，从$\egcd$算法出发。关键点在于要思考，什么是除法？]
    
 %%% 20250509, from a first course in CT.
  \item 设$\mathbf{C}$是一种向量空间，其向量长度为$10$（即每个向量有$10$个元素），且每个元素都是从$0$到$9$之间的整数。设矩阵$H$为
  \[
  H =
  \begin{bmatrix}
  	1 & 1 & 1 & 1 & 1 & 1 & 1 & 1 & 1 & 1 \\
  	1 & 2 & 3 & 4 & 5 & 6 & 7 & 8 & 9 & 10 
  \end{bmatrix}.
  \]
  设计$\mathbf{C}$的性质使得，任意$\mathbf{v} \in \mathbf{C}$当且仅当$H v \equiv \mathbf{0} \pmod{11}$。这样可保证，如果$\mathbf{C}$中向量在网络传输中发生了一个数位的差错，有一种巧妙的方法可以发现这种差错并纠正。请找到这种方法，并纠正$0617960587$这个向量发生的错误。
\end{problemset}
